% Abstract
Fine-grained execution feedback improves RL fine-tuning of code LLMs by providing token-level credit assignment, but this benefit depends on accurate error localization. We show that Python exception types partition into categories with dramatically different traceback reliability: U\_line errors (SyntaxError, NameError, etc.) have 100\% localization accuracy, while U\_ignore errors (RuntimeError, RecursionError, etc.) have only 20\% accuracy. Applying token-level penalties at unreliable locations injects gradient noise---we measure significantly lower gradient concentration at ground-truth bug locations for U\_ignore errors (1.398 vs 1.594, $p<10^{-13}$). We propose error-type gating: apply fine-grained feedback only for reliably-localized errors, falling back to coarse feedback otherwise. In proof-of-concept experiments on APPS with CodeT5-small, gating improves signal-to-noise ratio by 4.61\% and enables the model to reach 30\% pass@1 when the ungated baseline does not. Our work frames feedback granularity as a credit assignment reliability decision, providing both theoretical insight and a practical improvement for code RL training.
